\chapter{Zaključek}
\label{sec:zakljucek}

Podatki sledilnika pogleda vsebujejo časovne, prostorske in fiksacijske odnose med meritvami, ki jih klasična tabelarična struktura podatkov ne modelira neposredno.
Grafovske nevronske mreže omogočajo eksplicitno modeliranje relacij med meritvami, zato so naraven kandidat za obravnavo takšnih podatkov.
V nalogi smo preverili uporabnost grafovske predstavitve podatkov sledilnika pogleda pri binarni klasifikaciji valence.
Časovna okna meritev smo predstavili kot časovno-prostorske grafe in jih ovrednotili s štirimi grafovskimi modeli, ki se postopno nadgrajujejo od preproste proti kompleksnejšim različicam.
Grafovske modele smo primerjali s še tremi družinami modelov: izhodiščnima modeloma (naključni in večinski klasifikator), klasičnimi modeli strojnega učenja (\ModelName{SVM}, \ModelName{LightGBM}, \ModelName{MLP}) ter prednaučenima kodirnikoma signalov (\ModelName{GazeMAE}, \ModelName{MOMENT}).
Izvedli smo primerjalno evalvacijo vseh modelov na štirih delno prekrivajočih se množicah signalov: \emph{samo pogled}, \emph{samo zenici}, \emph{pogled in zenici} ter \emph{vsi signali}.

Glavni doprinosi naloge izhajajo direktno iz raziskovalnih vprašanj in so:

\begin{enumerate}
    \item časovno-prostorska grafovska predstavitev časovnih oken podatkov sledilnika pogleda;
    \item ovrednotenje arhitekturne lestvice grafovskih modelov za binarno klasifikacijo valence iz podatkov sledilnika pogleda;
    \item primerjava klasifikacijske uspešnosti med štirimi množicami signalov sledilnika pogleda;
    \item primerjava grafovskih modelov z negrafovskimi osnovnimi modeli in zamrznjenimi prednaučenimi predstavitvenimi modeli.
\end{enumerate}

Rezultati kažejo, da grafovski modeli iz meritev sledilnika pogleda izluščijo uporaben signal in presegajo izhodiščna klasifikatorja, njihova uspešnost pa je primerljiva z uspešnostjo zamrznjenih prednaučenih kodirnikov.
Kljub temu so se pri obravnavanem problemu najbolje izkazali klasični modeli strojnega učenja z ročno oblikovanimi značilkami.
Med preizkušenimi množicami signalov so bile najinformativnejše koordinate pogleda, velikosti zenic so prispevale nekaj dodatnih informacij, drugi razpoložljivi signali pa niso pokazali izrazite dodane vrednosti.

Grafovsko modeliranje podatkov sledilnika pogleda je še v veliki meri neraziskana, vendar smiselna nadaljnja raziskovalna smer.
V nadaljnjem delu bomo preverili drugačna pravila povezovanja, na primer povezovanje vozlišč s podobno dinamiko pogleda ali povezovanje meritev iste fiksacije z metavozliščem, s čimer bi graf razširili v tretjo dimenzijo.
Grafovsko predstavitev bi bilo mogoče dopolniti tudi z mežiki in sakadami ter preizkusiti dodatne grafovske arhitekture, predvsem grafovske transformerje.
Posebej zanimivo je samonadzorovano predtreniranje na signalih sledilnika pogleda, ki bi lahko zmanjšalo odvisnost od majhnih označenih podatkovnih zbirk in omogočilo uporabo neoznačenih podatkov.
Rezultate bi bilo treba preveriti še na več podatkovnih zbirkah, pri prenosu med nalogami in v pogojih uporabe v realnem času.

V nalogi smo tako pokazali, da je grafovska predstavitev smiselna za podatke sledilnika pogleda, saj je najboljši GNN pri klasifikaciji binarne valence dosegel točnost $65{,}9\,\%$.
